\summarychapter{14}{Peter as ``Pontifex Maximus''}{The Beginnings of the Papacy}

\begin{inbrief}
No institution has been so loved and so hated as the papacy. The chapter separates the early honor of
the Roman church from the later authority of its bishop, shows how councils, patriarchates and the
emperor's move to Constantinople prepared the ground, and presents Leo I, who gave the papal claim its
biblical basis and twice stood between Rome and the barbarians.
\question{On what does the pope's office rest, and at what point in history was that basis established?}
\end{inbrief}

In June 452 Attila the Hun, the ``scourge of God,'' crossed the Alps and advanced on Rome. At the Po
he received an embassy and found in it Bishop Leo, acting for an emperor who did nothing to defend the
old capital. Disease and famine in his army had probably already decided Attila to turn back, but no
one knew that; he granted the bishop's plea, spared Rome and left Italy. The bishop of Rome had taken
a new role and staked a claim on the future. Some revere the pope as Vicar of Christ and others
denounce him as Antichrist, but all agree that Leo marks a stage in an institution that has adapted
to the Roman Empire, the Germanic kingdoms, the modern nations and now the developing world.

\ssection{The Roman Catholic Claim}
The Roman Catholic Church teaches, as defined at the First Vatican Council (1870), that Christ founded
the papacy in Peter and that the bishop of Rome, his successor, holds supreme authority (primacy) over
the whole church. Eastern Orthodox and Protestants deny both points. Shelley's concern is history, and
history shows that papal rule was established by slow and painful stages; because Leo first supplied
its biblical and theological basis, it is misleading to speak of the papacy before him. The title
\emph{papa}, ``father,'' was used of any bishop and was reserved for Rome's only in the sixth century.

Honor must be distinguished from authority. The Roman church was honored early because it was the
church of the capital, the largest and richest in the West, with a name for orthodoxy and charity;
because of its size (by the mid-third century some 30,000 members, 150 clergy and 1,500 widows and
poor); and because writers from Irenaeus onward named Peter and Paul as its founders, a bishops' list
reaching back to them being a safeguard against gnostic secret traditions. Yet Irenaeus and Cyprian
disagreed with Rome when they thought it wrong, and before Constantine there is no conclusive evidence
that its bishop exercised jurisdiction outside Rome: honor, yes; jurisdiction, no.

Rome's rise was part of a growing church structure. Provincial councils of bishops met yearly from
the third century; in theory all bishops were equal, in practice those of apostolic churches and
large cities led. The church copied the map of the empire: the provincial town became the bishop's
town, the bishop of the metropolis supervised the province, and the great regions looked to Rome,
Carthage, Alexandria or Antioch. Country churches were served by priests (from \emph{presbyter}),
consecrated and controlled by the city bishop. General councils began after Constantine: Arles in 314
for the West, Nicaea in 325 for the whole church. At Nicaea the sees of Alexandria, Antioch and Rome were acknowledged as first in their own regions, with an honorary primacy for Jerusalem. The system of
patriarchates was thus confirmed.

\ssection{The Imperial Move East}
In 330 Constantine moved his residence to Constantinople, the New Rome. Political weight shifted
east, and the bishop of the new capital gained standing. Theodosius made Christianity the state
church, forbade sacrifices and closed the temples. In May 381 he opened a council at Constantinople to
which only eastern bishops were invited; Damasus of Rome sent no one. It renewed the Nicene Creed and
decreed that the bishop of Constantinople ranks immediately after the bishop of Rome ``because his
city is the New Rome.'' The West saw a political move. Damasus objected that a church's rank does not
depend on a city's status or a council's decree, and a Roman synod the next year declared that the
Roman church holds the primacy by the Lord's words, ``Thou art Peter, and upon this rock I will build
my church''---the first mention of the ``primacy of the Roman church.''

From then on the two churches went different ways. Constantinople relied on its political position
and its patriarch grew less independent as religion and politics intertwined. In the West the
weakness of the empire left the bishop of Rome without a rival and increasingly independent; unable
to argue from Rome the imperial city, he argued from Scripture and tradition, the primacy of Peter.
Damasus (366--384) fused Roman civic pride with Christianity and pointed to Constantine's basilicas
over the graves of Peter and Paul: by their martyrdom Rome had the better right to claim them as
citizens.

\ssection{Leo's Case for Primacy}
Leo, a nobleman from north of Rome, was on an imperial mission in Gaul when the Roman clergy elected
him. In his first sermon he praised the glory of Peter, ``in whose chair his power lives on.''
Appealing to three gospel texts (Matthew 16, Luke 22, John 21), he laid the theoretical foundation of
primacy: Christ founded the church on Peter as its lasting rock, and Peter's authority passes to each bishop of Rome. Where Constantine had used the church as a tool, Leo raised the Roman bishop's office as
far, in theory, as it could go. Shelley notes three difficulties. The Gospels say that leadership
among Christ's followers is humble service, not lordship. Peter remained unstable, rebuked as
``Satan,'' denying his Lord and criticized by Paul. And the texts nowhere tie the authority to an
office, the bishopric of Rome.

The argument nevertheless suited the hour. The court at Ravenna, pressed by the barbarians, needed
any authority that could hold the West together, and in 445 the emperor Valentinian III decreed that
bishops who would not come to the papal court be compelled, grounding the primacy of the Apostolic See
on Peter, on the dignity of Rome and on an unnamed ``Holy Synod.'' Leo's claim became law. At
Chalcedon in October 451 his letters dominated, and the 350 bishops cried, ``Peter has spoken through
Leo.'' Yet on October 30 the council also ruled that Constantinople's bishop, since his city was the New Rome, should have the same authority as Rome's. Leo's legate protested in vain. Christianity now had two heads, the Roman
church of the West and the Greek church of the East.

\ssection{Standing Amidst the Ruins}
Three years after Attila came the Vandals, a Scandinavian people driven across Gaul and Spain and
settled in North Africa. At the end of March 455 their king Gaiseric sailed with a hundred ships and
landed north of the Tiber. Rome panicked, the troops mutinied, and the emperor Maximus was killed by
his own bodyguard while fleeing. On June 2 the Vandals entered without resistance. Leo, at the head
of his priests, met Gaiseric at the gate and begged him to restrain his men and not to burn the city.
The king answered, ``Fourteen days' looting!'' The Vandals stripped the palaces and temples, took the
gilded roof of the Capitol and the vessels of Solomon's temple, and carried the empress, her daughters
and senators off to Carthage for ransom. But the city was not burned, there was no massacre, and only
a few churches were plundered.

At the thanksgiving service few came. Leo rebuked the absent for valuing ``the devils higher than the
Apostles'' and for crediting their deliverance to the stars rather than to the mercy of God, who had
softened the barbarians. He did not mention himself and did not need to. He had taken over the old pagan title Pontifex Maximus, once borne by the chief priest of the empire's religion, and its meaning escaped no one: the Eternal City was now in the care of Leo and not of the emperor. Peter, Shelley says, had come to power.

\begin{keyfacts}
\keypeople{Leo I (the Great); Attila the Hun; Damasus; Theodosius; Valentinian III; Gaiseric, king of
the Vandals; the emperor Maximus; Irenaeus; Cyprian.}
\keydates{%
  \kd{314}{Council of Arles, a general council of the western churches}
  \kd{325}{Nicaea recognizes the leading sees of Alexandria, Antioch and Rome}
  \kd{381}{Council of Constantinople ranks its bishop next after Rome's}
  \kd{382}{A Roman synod first speaks of the ``primacy of the Roman church''}
  \kd{440--461}{Leo I bishop of Rome}
  \kd{445}{Valentinian III makes the primacy of the Roman see imperial law}
  \kd{451}{Chalcedon: ``Peter has spoken through Leo''; Constantinople made equal to Rome}
  \kd{452}{Leo persuades Attila to spare Rome}
  \kd{455}{The Vandals plunder Rome for fourteen days; Leo prevents burning and massacre}
  \kd{1870}{First Vatican Council defines papal primacy}}
\keyterms{Papacy; primacy; \emph{papa}; honor and jurisdiction; patriarchates; general council;
New Rome; ``Thou art Peter''; Pontifex Maximus.}
\end{keyfacts}
